% GENERATED FILE. Edit canonical lessons, then run npm run generate:books.
% canonical-source-hash: fnv1a64:84d015f4c874bc94
% canonical-lessons: ML-C491-mal, ML-C491-ophar, ML-C491-adayavilppana, ML-C491-vilppanakkaran, ML-C491-ittunokkuka

\chapter{At the Mall}
\label{ch:ml-at-the-mall}
\hlchaptermodality{491}

\begin{chapteropening}
\textbf{By the end of this chapter:} \emph{I can say a shopping mall, a special offer, a clearance sale, a salesman and to try on.}

Complete the last lesson of chapter 491: I can say a shopping mall, a special offer, a clearance sale, a salesman and to try on.
\end{chapteropening}

\section[\texorpdfstring{\ml{മാൾ} (māḷ)}{മാൾ (māḷ)}]{\ml{മാൾ} (māḷ) — a shopping mall}

\label{lesson:ML-C491-mal}



\begin{quote}
Before the new one: say the Malayalam for to send back, then the Malayalam for to contact.
\end{quote}

\subsection*{You'll want to know: \ml{മാൾ}}
\textbf{\ml{മാൾ}} — \emph{māḷ} — \textquotedblleft{}a shopping mall\textquotedblright{}.

\begin{quote}
\ml{ഞായറാഴ്ച} \ml{മാളിൽ} \ml{വലിയ} \ml{തിരക്കായിരുന്നു}.
\end{quote}

\begin{quote}
\emph{ñāyaṟāḻca māḷil valiya tirakkāyirunnu.}
\end{quote}

\begin{quote}
The mall was very crowded on Sunday.
\end{quote}

A loan ending in a chillu: \ml{മാൾ}, \ml{മാളിൽ}, in the mall.

\begin{cousinweb}
\textbf{\ml{മാൾ}} (\emph{māḷ}) was borrowed from English \emph{mall}, a large building full of shops.
\end{cousinweb}

\begin{grammarlens}[title={What you've built}]
One more word for getting around.
\end{grammarlens}

\subsection*{Your turn}
\begin{itemize}
\raggedright
  \item \emph{Say it:} \emph{māḷ}
  \item \emph{Say it:} \emph{māḷ}, once more
  \item \emph{Say it:} \emph{ñāyaṟāḻca māḷil valiya tirakkāyirunnu.} — The mall was very crowded on Sunday.
  \item \emph{Recall:} say the Malayalam for to send back, then the Malayalam for to contact, then say \emph{māḷ} again
\end{itemize}

\begin{tcolorbox}[breakable,colback=teal!4,colframe=teal!35!black,title={Before you move on}]
What does \ml{മാൾ} mean? (\textquotedblleft{}A shopping mall\textquotedblright{}.) Say it once more.
\end{tcolorbox}
\section[\texorpdfstring{\ml{ഓഫർ} (ōphar)}{ഓഫർ (ōphar)}]{\ml{ഓഫർ} (ōphar) — a special offer}

\label{lesson:ML-C491-ophar}



\begin{quote}
Before the new one: say the Malayalam for to contact, then the Malayalam for a shopping mall.
\end{quote}

\subsection*{You'll want to know: \ml{ഓഫർ}}
\textbf{\ml{ഓഫർ}} — \emph{ōphar} — \textquotedblleft{}a special offer\textquotedblright{}.

\begin{quote}
\ml{ആ} \ml{കടയിൽ} \ml{ഒന്ന്} \ml{വാങ്ങിയാൽ} \ml{ഒന്ന്} \ml{വെറുതെ} \ml{എന്ന} \ml{ഓഫർ} \ml{ഉണ്ട്}.
\end{quote}

\begin{quote}
\emph{ā kaṭayil onnŭ vāṅṅiyāl onnŭ veṟute enna ōphar uṇṭŭ.}
\end{quote}

\begin{quote}
That shop has an offer: buy one and get one free.
\end{quote}

Shops put up an \ml{ഓഫർ} board at festival time.

\begin{cousinweb}
\textbf{\ml{ഓഫർ}} (\emph{ōphar}) was borrowed from English \emph{offer}, a price cut for a short time.
\end{cousinweb}

\begin{grammarlens}[title={What you've built}]
One more word for getting around.
\end{grammarlens}

\subsection*{Your turn}
\begin{itemize}
\raggedright
  \item \emph{Say it:} \emph{ōphar}
  \item \emph{Say it:} \emph{ōphar}, once more
  \item \emph{Say it:} \emph{ā kaṭayil onnŭ vāṅṅiyāl onnŭ veṟute enna ōphar uṇṭŭ.} — That shop has an offer: buy one and get one free.
  \item \emph{Recall:} say the Malayalam for to contact, then the Malayalam for a shopping mall, then say \emph{ōphar} again
\end{itemize}

\begin{tcolorbox}[breakable,colback=teal!4,colframe=teal!35!black,title={Before you move on}]
What does \ml{ഓഫർ} mean? (\textquotedblleft{}A special offer\textquotedblright{}.) Say it once more.
\end{tcolorbox}
\section[\texorpdfstring{\ml{ആദായവിൽപ്പന} (ādāyavilppana)}{ആദായവിൽപ്പന (ādāyavilppana)}]{\ml{ആദായവിൽപ്പന} (ādāyavilppana) — a clearance sale}

\label{lesson:ML-C491-adayavilppana}



\begin{quote}
Before the new one: say the Malayalam for a shopping mall, then the Malayalam for a special offer.
\end{quote}

\subsection*{You'll want to know: \ml{ആദായവിൽപ്പന}}
\textbf{\ml{ആദായവിൽപ്പന}} — \emph{ādāyavilppana} — \textquotedblleft{}a clearance sale\textquotedblright{}.

\begin{quote}
\ml{അടുത്ത} \ml{മാസം} \ml{തുണിക്കടകളിൽ} \ml{ആദായവിൽപ്പന} \ml{ഉണ്ടാകും}.
\end{quote}

\begin{quote}
\emph{aṭutta māsaṁ tuṇikkaṭakaḷil ādāyavilppana uṇṭākuṁ.}
\end{quote}

\begin{quote}
Next month the clothes shops will have clearance sales.
\end{quote}

\ml{ആദായം}, gain, plus \ml{വിൽപ്പന}, a sale: a sale where the buyer gains.

\begin{cousinweb}
\textbf{\ml{ആദായവിൽപ്പന}} (\emph{ādāyavilppana}) is a Malayalam compound built from \ml{ആദായം}, profit, from Sanskrit \emph{ādāya}, \textquotedblleft{}taking\textquotedblright{} (Monier-Williams), and the native word \ml{വിൽപ്പന}, a sale (DEDR 5421).
\end{cousinweb}

\begin{grammarlens}[title={What you've built}]
One more word for getting around.
\end{grammarlens}

\subsection*{Your turn}
\begin{itemize}
\raggedright
  \item \emph{Say it:} \emph{ādāyavilppana}
  \item \emph{Say it:} \emph{ādāyavilppana}, once more
  \item \emph{Say it:} \emph{aṭutta māsaṁ tuṇikkaṭakaḷil ādāyavilppana uṇṭākuṁ.} — Next month the clothes shops will have clearance sales.
  \item \emph{Recall:} say the Malayalam for a shopping mall, then the Malayalam for a special offer, then say \emph{ādāyavilppana} again
\end{itemize}

\begin{tcolorbox}[breakable,colback=teal!4,colframe=teal!35!black,title={Before you move on}]
What does \ml{ആദായവിൽപ്പന} mean? (\textquotedblleft{}A clearance sale\textquotedblright{}.) Say it once more.
\end{tcolorbox}
\section[\texorpdfstring{\ml{വിൽപ്പനക്കാരൻ} (vilppanakkāran)}{വിൽപ്പനക്കാരൻ (vilppanakkāran)}]{\ml{വിൽപ്പനക്കാരൻ} (vilppanakkāran) — a salesman}

\label{lesson:ML-C491-vilppanakkaran}



\begin{quote}
Before the new one: say the Malayalam for a special offer, then the Malayalam for a clearance sale.
\end{quote}

\subsection*{You'll want to know: \ml{വിൽപ്പനക്കാരൻ}}
\textbf{\ml{വിൽപ്പനക്കാരൻ}} — \emph{vilppanakkāran} — \textquotedblleft{}a salesman\textquotedblright{}.

\begin{quote}
\ml{വിൽപ്പനക്കാരൻ} \ml{എനിക്ക്} \ml{പല} \ml{നിറങ്ങളിലുള്ള} \ml{ഷർട്ടുകൾ} \ml{കാണിച്ചു}.
\end{quote}

\begin{quote}
\emph{vilppanakkāran enikkŭ pala niṟaṅṅaḷiluḷḷa ṣarṭṭukaḷ kāṇiccu.}
\end{quote}

\begin{quote}
The salesman showed me shirts in many colours.
\end{quote}

\ml{വിൽപ്പന}, a sale, plus -\ml{കാരൻ}, one who: a saleswoman is a \ml{വിൽപ്പനക്കാരി}.

\begin{cousinweb}
\textbf{\ml{വിൽപ്പനക്കാരൻ}} (\emph{vilppanakkāran}) is native Dravidian, built from \ml{വിൽപ്പന}, a sale (DEDR 5421), with -\ml{കാരൻ}, one who does it.
\end{cousinweb}

\begin{grammarlens}[title={What you've built}]
One more word for getting around.
\end{grammarlens}

\subsection*{Your turn}
\begin{itemize}
\raggedright
  \item \emph{Say it:} \emph{vilppanakkāran}
  \item \emph{Say it:} \emph{vilppanakkāran}, once more
  \item \emph{Say it:} \emph{vilppanakkāran enikkŭ pala niṟaṅṅaḷiluḷḷa ṣarṭṭukaḷ kāṇiccu.} — The salesman showed me shirts in many colours.
  \item \emph{Recall:} say the Malayalam for a special offer, then the Malayalam for a clearance sale, then say \emph{vilppanakkāran} again
\end{itemize}

\begin{tcolorbox}[breakable,colback=teal!4,colframe=teal!35!black,title={Before you move on}]
What does \ml{വിൽപ്പനക്കാരൻ} mean? (\textquotedblleft{}A salesman\textquotedblright{}.) Say it once more.
\end{tcolorbox}
\section[\texorpdfstring{\ml{ഇട്ടുനോക്കുക} (iṭṭunōkkuka)}{ഇട്ടുനോക്കുക (iṭṭunōkkuka)}]{\ml{ഇട്ടുനോക്കുക} (iṭṭunōkkuka) — to try on}

\label{lesson:ML-C491-ittunokkuka}



\begin{quote}
Before the new one: say the Malayalam for a clearance sale, then the Malayalam for a salesman.
\end{quote}

\subsection*{You'll want to know: \ml{ഇട്ടുനോക്കുക}}
\textbf{\ml{ഇട്ടുനോക്കുക}} — \emph{iṭṭunōkkuka} — \textquotedblleft{}to try on\textquotedblright{}.

\begin{quote}
\ml{വാങ്ങുന്നതിന്} \ml{മുൻപ്} \ml{ഞാൻ} \ml{ആ} \ml{ഷർട്ട്} \ml{ഇട്ടുനോക്കി}.
\end{quote}

\begin{quote}
\emph{vāṅṅunnatinŭ munpŭ ñān ā ṣarṭṭŭ iṭṭunōkki.}
\end{quote}

\begin{quote}
I tried the shirt on before buying it.
\end{quote}

Past: \ml{ഇട്ടുനോക്കി}, \emph{iṭṭunōkki}.

Literally to put on and look: \ml{ഇട്ട്}, having put on, plus \ml{നോക്കുക}. Any verb plus \ml{നോക്കുക} means to try doing it.

\begin{cousinweb}
\textbf{\ml{ഇട്ടുനോക്കുക}} (\emph{iṭṭunōkkuka}) is a native Dravidian compound of \ml{ഇടുക}, to put on (DEDR 442), and \ml{നോക്കുക}, to look (DEDR 3794).
\end{cousinweb}

\begin{grammarlens}[title={What you've built}]
One more word for getting around.
\end{grammarlens}

\subsection*{Your turn}
\begin{itemize}
\raggedright
  \item \emph{Say it:} \emph{iṭṭunōkkuka}
  \item \emph{Say it:} \emph{iṭṭunōkkuka}, once more
  \item \emph{Say it:} \emph{vāṅṅunnatinŭ munpŭ ñān ā ṣarṭṭŭ iṭṭunōkki.} — I tried the shirt on before buying it.
  \item \emph{Recall:} say the Malayalam for a clearance sale, then the Malayalam for a salesman, then say \emph{iṭṭunōkkuka} again
\end{itemize}

\begin{tcolorbox}[breakable,colback=teal!4,colframe=teal!35!black,title={Before you move on}]
What does \ml{ഇട്ടുനോക്കുക} mean? (\textquotedblleft{}To try on\textquotedblright{}.) Say it once more.
\end{tcolorbox}
